% =====================================================================
% Complete formulation and proofs. Included as the appendix of main.tex
% and as the body of supplement.tex. Every equation matches the code in
% genco_rss/ (file and function given in brackets).
% =====================================================================

\section{DC Model of the Robust Secure Set}
\label{app:dc}

\subsection{Sets, data and decision variables}
Buses $\mathcal N$ ($|\mathcal N|=118$) with reference bus $n_0$; branches $\mathcal L$
($|\mathcal L|=186$) with reactance $x_\ell$ and MVA rating $r_\ell$ (p.u., base 100\,MVA);
dispatchable units $\mathcal G$ ($|\mathcal G|=19$) with the balancing (slack) unit
$g_0\in\mathcal G$ located at $n_0$; PV buses $\mathcal R$ ($|\mathcal R|=20$); load buses
$\mathcal D$; topologies $s\in\mathcal S=\{0,1,2\}$ (base case; one circuit of 69--70 lost;
one circuit of 100--103 lost). Incidence matrices $C_G\in\{0,1\}^{|\mathcal N|\times|\mathcal G|}$,
$C_R$, $C_D$ map units, PV injections and loads to buses.

Uncertain vector (nature) $p=[P_R;P_D]\in\mathbb R^{n_p}$, $n_p=|\mathcal R|+|\mathcal D|$.
Dispatch (operator) $u\in\mathbb R^{n_u}_{\ge 0}$, $n_u=|\mathcal G|$
(plus $2|\mathcal R|$ battery variables in \apref{app:bess}).

\subsection{Power transfer distribution factors}
[\texttt{grid.ptdf}] Let $C\in\{-1,0,1\}^{|\mathcal L|\times|\mathcal N|}$ be the
branch--bus incidence matrix and $X_s=\mathrm{diag}(x^s_\ell)$ with
\begin{equation}
x^s_\ell=\begin{cases}2x_\ell & \ell \text{ is the corridor of } s\ (\text{one of two identical parallel circuits lost}),\\
x_\ell & \text{otherwise,}\end{cases}
\label{eq:xs}
\end{equation}
since two identical circuits in parallel have half the reactance of one. As in the
dissertation data, the corridor rating is kept.
With $B_f^s=X_s^{-1}C$, $B_{\rm bus}^s=C^\top X_s^{-1}C$ and $\mathcal N'=\mathcal N\setminus\{n_0\}$,
\begin{equation}
H_s[:,\mathcal N']=B_f^s[:,\mathcal N']\big(B^s_{\rm bus}[\mathcal N',\mathcal N']\big)^{-1},\qquad H_s[:,n_0]=0 .
\label{eq:ptdf}
\end{equation}
The DC branch flows of an operating point in topology $s$ are
\begin{equation}
f_s(p,u)=H_s\big(C_Gu+C_RP_R-C_DP_D\big).
\label{eq:flow}
\end{equation}
Because the slack unit sits at $n_0$, column $H_sC_G[:,g_0]=0$: the slack unit does not
appear in any line constraint (used in Lemma~\ref{lem:robust}).

\subsection{Constraints}
[\texttt{grid.rss\_matrices}] For derating factors $\kappa_s\in(0,1]^{|\mathcal L|}$
($\kappa\equiv1$ in the DC study, \apref{app:ac} otherwise):
\begin{align}
&-\kappa_s\circ r\le f_s(p,u)\le\kappa_s\circ r, && s\in\mathcal S, \label{eq:line}\\
&\mathbf 1^\top u+\mathbf 1^\top P_R-\mathbf 1^\top P_D=0, && \text{(lossless balance)} \label{eq:bal}\\
&\underline u\le u\le\overline u. && \label{eq:genb}
\end{align}
[\texttt{grid.gen\_limits}, \texttt{rss.dc\_opf}] With $u^0$ the linear-cost DC-OPF dispatch at
the box centre and ramp fraction $\rho=0.3$:
\begin{equation}
\underline u_k=\max(P^{\min}_k,u^0_k-\rho P^{\max}_k),\quad
\overline u_k=\min(P^{\max}_k,u^0_k+\rho P^{\max}_k)\ (k\neq g_0),\quad
[\underline u_{g_0},\overline u_{g_0}]=[P^{\min}_{g_0},P^{\max}_{g_0}].
\label{eq:ramp}
\end{equation}
Stacking (\ref{eq:line})--(\ref{eq:genb}) row-wise gives the polyhedral system
\begin{equation}
Ap+Bu\le b,\qquad
A=\begin{bmatrix}A_\ell\\-A_\ell\\ a^\top\\-a^\top\\0\\0\end{bmatrix},\
B=\begin{bmatrix}B_\ell\\-B_\ell\\ \mathbf 1^\top\\-\mathbf 1^\top\\-I\\I\end{bmatrix},\
b=\begin{bmatrix}\kappa\circ r\\ \kappa\circ r\\0\\0\\-\underline u\\ \overline u\end{bmatrix},
\label{eq:Abb}
\end{equation}
with $A_\ell=[H_sC_R,\,-H_sC_D]_{s\in\mathcal S}$, $B_\ell=[H_sC_G]_{s\in\mathcal S}$
(stacked over $s$), $a=[\mathbf 1;-\mathbf 1]$. The follower's feasible set is
\begin{equation}
\mathcal U(p)=\{u\in\mathbb R^{n_u}_{\ge0}: Bu\le b-Ap\}.
\label{eq:U}
\end{equation}

\subsection{Uncertainty set}
[\texttt{grid.scenario\_box}] For hour $h$, with PV forecast $\hat P_{R}$, capacity
$\bar P_R$, forecast load $\hat P_D$, $e_{\rm PV}(h)=0.3\,h/24$ and $e_D(h)=0.03\,h/24$:
\begin{equation}
\mathcal P_h=\Big\{p:\ \big[\hat P_R-e_{\rm PV}\bar P_R\big]_0^{\bar P_R}\le P_R\le\big[\hat P_R+e_{\rm PV}\bar P_R\big]_0^{\bar P_R},\
(1-e_D)\hat P_D\le P_D\le(1+e_D)\hat P_D\Big\},
\label{eq:box}
\end{equation}
where $[\cdot]_0^{\bar P}$ clips to $[0,\bar P]$ and $P_R\equiv0$ for $h\notin[7,18]$.
$\mathcal P_h=[\underline p,\overline p]$ is a box with vertex set $V(\mathcal P_h)$.

\subsection{Robust secure set}
The RSS index is the total output of the non-balancing units,
\begin{equation}
\alpha=d^\top u,\qquad d_k=1\ (k\neq g_0),\ d_{g_0}=0 .
\label{eq:index}
\end{equation}
\begin{assumption}\label{as:feas}
$\mathcal U(p)\neq\emptyset$ for every $p\in\mathcal P_h$ (the follower can always serve the
load; checked by Benders, which otherwise reports infeasibility).
\end{assumption}
Define the follower value functions and the RSS limits
\begin{align}
\varphi(p)&=\max\{d^\top u:u\in\mathcal U(p)\}, & \psi(p)&=\min\{d^\top u:u\in\mathcal U(p)\},\label{eq:vf}\\
\overline\alpha_h&=\min_{p\in\mathcal P_h}\varphi(p), & \underline\alpha_h&=\max_{p\in\mathcal P_h}\psi(p). \label{eq:rssdef}
\end{align}

\begin{theorem}[Meaning of the RSS interval]\label{thm:rss}
Under Assumption~\ref{as:feas}, a value $\alpha$ is achievable for \emph{every}
realisation, i.e.\ $\forall p\in\mathcal P_h\ \exists u\in\mathcal U(p): d^\top u=\alpha$, if
and only if $\alpha\in[\underline\alpha_h,\overline\alpha_h]$. In particular the robust
secure set is empty if and only if $\underline\alpha_h>\overline\alpha_h$.
\end{theorem}
\begin{proof}
Fix $p$. $\mathcal U(p)$ is a non-empty polytope (bounded by (\ref{eq:genb})), hence
compact and convex; its image under the linear map $u\mapsto d^\top u$ is a compact
interval, which by (\ref{eq:vf}) equals $[\psi(p),\varphi(p)]$.
($\Leftarrow$) If $\underline\alpha_h\le\alpha\le\overline\alpha_h$, then for every $p$,
$\psi(p)\le\underline\alpha_h\le\alpha\le\overline\alpha_h\le\varphi(p)$, so $\alpha$ lies in the
image interval and some $u\in\mathcal U(p)$ attains it.
($\Rightarrow$) If $\alpha>\overline\alpha_h$, take $p^\star\in\arg\min\varphi$ (exists by
Lemma~\ref{lem:vertex}); then $\alpha>\varphi(p^\star)$ and no $u\in\mathcal U(p^\star)$ attains
$\alpha$. The case $\alpha<\underline\alpha_h$ is symmetric with $\arg\max\psi$.
\end{proof}

\begin{lemma}[Worst cases are vertices]\label{lem:vertex}
$\varphi$ is concave and $\psi$ is convex on $\mathcal P_h$. Hence
\[
\overline\alpha_h=\min_{p\in V(\mathcal P_h)}\varphi(p),\qquad
\underline\alpha_h=\max_{p\in V(\mathcal P_h)}\psi(p).
\]
\end{lemma}
\begin{proof}
Let $p_1,p_2\in\mathcal P_h$, $\theta\in[0,1]$, and $u_i\in\mathcal U(p_i)$ optimal for
$\varphi(p_i)$. Then $u_\theta=\theta u_1+(1-\theta)u_2\ge0$ and
$Bu_\theta=\theta Bu_1+(1-\theta)Bu_2\le b-A(\theta p_1+(1-\theta)p_2)$, so
$u_\theta\in\mathcal U(p_\theta)$ and $\varphi(p_\theta)\ge d^\top u_\theta=\theta\varphi(p_1)+(1-\theta)\varphi(p_2)$:
$\varphi$ is concave. The same argument with minimisers shows that $\psi$ is convex.
A concave function on a polytope attains its minimum at a vertex (write any point as a
convex combination of vertices and use concavity: $\varphi(\sum\theta_jv_j)\ge\sum\theta_j\varphi(v_j)\ge\min_j\varphi(v_j)$);
symmetrically a convex function attains its maximum at a vertex.
\end{proof}

\section{Bilevel Form, KKT Conditions and the Big-M Baseline}
\label{app:bilevel}
Both limits are instances of the linear bilevel program [\texttt{benders.solve}]
\begin{equation}
\mathrm{OPT}=\max_{p\in\mathcal P_h,\,u}\ c_1^\top p+d_1^\top u\quad\text{s.t.}\quad
u\in\arg\max_{v\in\mathcal U(p)} d_2^\top v,
\label{eq:bilevel}
\end{equation}
with $c_1=0$ and $(d_1,d_2)=(-d,d)$ for the upper limit ($\overline\alpha_h=-\mathrm{OPT}$) and
$(d_1,d_2)=(d,-d)$ for the lower limit ($\underline\alpha_h=\mathrm{OPT}$).
\begin{remark}\label{rem:optimistic}
Since $d_1=-d_2$, every follower optimum gives the leader the same value
$d_1^\top u=-\max_{v}d_2^\top v$; optimistic and pessimistic bilevel solutions coincide. For
general data we use the optimistic convention: the attained value at $p$ is
\begin{equation}
F(p)=c_1^\top p+\max\{d_1^\top u:\ u\in\mathcal U(p),\ d_2^\top u=\max_{v\in\mathcal U(p)}d_2^\top v\},
\label{eq:Fp}
\end{equation}
computed by two LPs [\texttt{benders.follower\_response}].
\end{remark}
\textbf{Follower dual and KKT.} The follower LP $\max\{d_2^\top u: Bu\le b-Ap,\ u\ge0\}$ has the dual
$\min\{\lambda^\top(b-Ap):\ B^\top\lambda\ge d_2,\ \lambda\ge0\}$. A feasible $u$ is optimal iff
there is $\lambda$ with
\begin{equation}
\lambda\ge0,\quad B^\top\lambda\ge d_2,\quad u_k\,(B^\top\lambda-d_2)_k=0\ \forall k,\quad
\lambda_i\,(b-Ap-Bu)_i=0\ \forall i .
\label{eq:kkt}
\end{equation}
Let $\mathcal R=\{i: B_i\neq0\}$. Rows with $B_i=0$ do not enter $B^\top\lambda$; setting
$\lambda_i=0$ there preserves (\ref{eq:kkt}), so w.l.o.g.\ $\lambda$ is supported on $\mathcal R$.

\textbf{Big-M KKT MILP of the dissertation} [\texttt{rss.bigm\_milp}, \texttt{RSS\_MILP2.m}].
With $x=(p,u,\lambda)$, $G(x)=\big(0;\ B^\top\lambda-d_2;\ b-Ap-Bu\big)$ and binaries $w$:
\begin{equation}
\max\ c_1^\top p+d_1^\top u\ \ \text{s.t.}\ \ Ap+Bu\le b,\ \ 0\le x\le Lw,\ \ 0\le G(x)\le L(\mathbf 1-w),\ \ p\in\mathcal P_h,
\label{eq:bigm}
\end{equation}
$L=5\cdot10^4$. If $L$ bounds every component of some optimal KKT pair, (\ref{eq:bigm}) is exact.
\begin{remark}[Why the big-M MILP misreports]\label{rem:bigm}
(i) If $L$ is smaller than the components of every optimal KKT pair, (\ref{eq:bigm})
cuts the optimum off~\cite{pineda2019bigm}. (ii) A MILP solver accepts $w_j\le\varepsilon_I$ as
$0$ (integrality tolerance, $\varepsilon_I=10^{-6}$ in HiGHS); then $x_j\le L\varepsilon_I=0.05$
while $G_j(x)>0$ is allowed, i.e.\ complementarity can be violated by up to $L\varepsilon_I$
per pair. The reported $u$ is then not follower-optimal and the reported objective is
not attained. The reporting errors measured in E1 (up to 0.063\,p.u.) are of this order.
Our method never reports a value that is not attained: the incumbent is always
$F(p)$ of (\ref{eq:Fp}), computed by LP.
\end{remark}

\section{Two-Phase Combinatorial Benders Decomposition}
\label{app:benders}

\subsection{Exact bounds}
For $i\in\mathcal R$ the largest possible slack over the box is
\begin{equation}
S_i=\max\Big(0,\ b_i-\min_{p\in\mathcal P_h,\,0\le u\le\overline u}(A_ip+B_iu)\Big)
=\max\Big(0,\ b_i-\textstyle\sum_j\min(A_{ij}\underline p_j,A_{ij}\overline p_j)-\sum_k\min(0,B_{ik}\overline u_k)\Big),
\label{eq:S}
\end{equation}
and $\overline U_k=\overline u_k$. Both are computed from the data; no constant is chosen.

\subsection{Phase-1 master}
For a finite set $\mathcal C$ of cuts, the master problem $\mathrm{MP}(\mathcal C)$ has
variables $p,u$ and binaries $z^u\in\{0,1\}^{n_u}$, $z^\lambda\in\{0,1\}^{|\mathcal R|}$:
\begin{align}
\max\ & c_1^\top p+d_1^\top u \label{eq:mp-obj}\\
\text{s.t. } & Ap+Bu\le b,\quad p\in\mathcal P_h,\quad u\ge0, \tag{M1}\label{eq:M1}\\
& u_k\le\overline U_k\,z^u_k, && k=1..n_u \tag{M2}\label{eq:M2}\\
& b_i-A_ip-B_iu\le S_i\,(1-z^\lambda_i), && i\in\mathcal R \tag{M3}\label{eq:M3}\\
& z^\lambda_a+z^\lambda_j\le1, && (a,j)\in\Pi \tag{M4}\label{eq:M4}\\
& d_2^\top u\ge\varphi_R, && \text{if }\varphi_R>-\infty \tag{M5}\label{eq:M5}\\
& \textstyle\sum_{k\in C_u}(1-z^u_k)+\sum_{i\in C_\ell}z^\lambda_i\ge1, && (C_u,C_\ell)\in\mathcal C \tag{M6}\label{eq:M6}
\end{align}
where $\Pi$ is the set of opposite row pairs, $[A_a\ B_a]=-[A_j\ B_j]$ with $b_a+b_j>0$
(e.g.\ the two sides of a line limit), and $\varphi_R$ is defined in (\ref{eq:phiR}).
Interpretation: $z^u_k=0$ forces $u_k=0$; $z^\lambda_i=1$ forces row $i$ to be binding;
$z^\lambda_i=0$ will mean $\lambda_i=0$.

\subsection{Subproblem, Farkas certificate and cut}
For index sets $J_u\subseteq\{1..n_u\}$ (units allowed to be positive) and
$J_\ell\subseteq\mathcal R$ (rows with positive slack) define the dual system
\begin{equation}
\mathrm{DF}(J_u,J_\ell):\quad \lambda\in\mathbb R^{\mathcal R}_{\ge0},\quad \lambda_i=0\ (i\in J_\ell),\quad
B_{\mathcal R}^\top\lambda\ge d_2,\quad (B_{\mathcal R}^\top\lambda)_k=d_{2,k}\ (k\in J_u).
\label{eq:DF}
\end{equation}
[\texttt{benders.\_dual\_feasible}] Given a master solution $(\hat p,\hat u,\hat z)$ with slack
$\hat s=b-A\hat p-B\hat u$, the subproblem checks $\mathrm{DF}(J_u(\hat u),J_\ell(\hat s))$,
$J_u(\hat u)=\{k:\hat u_k>0\}$, $J_\ell(\hat s)=\{i\in\mathcal R:\hat s_i>0\}$.

\begin{lemma}[Farkas certificate]\label{lem:farkas}
Let $\tilde B=B_{\mathcal R\setminus J_\ell}$ (the rows whose multiplier may be positive).
$\mathrm{DF}(J_u,J_\ell)$ is infeasible if and only if there is $y\in\mathbb R^{n_u}$ with
\begin{equation}
\tilde By\le0,\qquad y_k\ge0\ (k\notin J_u),\qquad d_2^\top y>0 .
\label{eq:farkas}
\end{equation}
\end{lemma}
\begin{proof}
Eliminating $\lambda_{J_\ell}=0$ and introducing surplus $t\ge0$ with $t_{J_u}=0$,
$\mathrm{DF}$ reads $\{(\lambda,t_{\bar J_u})\ge0:\ \tilde B^\top\lambda-I_{:,\bar J_u}t_{\bar J_u}=d_2\}$,
$\bar J_u$ the complement of $J_u$. By Farkas' lemma, $\{x\ge0:Mx=d_2\}$ is infeasible iff
some $y$ has $M^\top y\le0$ and $d_2^\top y>0$; with $M=[\tilde B^\top,\,-I_{:,\bar J_u}]$ this is
(\ref{eq:farkas}).
\end{proof}
[\texttt{benders.\_farkas\_subset}] A certificate also identifies small infeasible subsystems:
an elastic LP relaxes each restriction $\lambda_i=0$ ($i\in J_\ell$) and
$(B^\top\lambda)_k=d_{2,k}$ ($k\in J_u$) by a non-negative variable and minimises their sum; the
restrictions with a positive dual form $C_\ell\subseteq J_\ell$, $C_u\subseteq J_u$. A second LP checks
that $\mathrm{DF}(C_u,C_\ell)$ is infeasible; if not, the full sets $C_u=J_u$, $C_\ell=J_\ell$ are used,
which is always correct. Smaller sets give stronger cuts.

If $\mathrm{DF}(C_u,C_\ell)$ is infeasible, the \emph{combinatorial Benders cut}
(\ref{eq:M6})~\cite{codato2006} is added. Equivalently
$\sum_{k\in C_u}z^u_k-\sum_{i\in C_\ell}z^\lambda_i\le|C_u|-1$.

\subsection{Robust value-function cut}
[\texttt{benders.robust\_bound}] Let $F\subseteq\{k:d_{2,k}=0\}$ be a set of balancing
columns such that every row of $B$ either has $B_{iF}=0$, or is an own bound row of
one $f\in F$ (with $A_i=0$), or belongs to the single equality pair $(i_1,i_2)$ with
$[A_{i_1}\ B_{i_1}]=-[A_{i_2}\ B_{i_2}]$ (the balance (\ref{eq:bal})). For the RSS,
$F=\{g_0\}$ (the slack unit has a zero PTDF column, \apref{app:dc}); if the
structure is not met, $F=\emptyset$. Let $[\underline w_f,\overline w_f]$ be the own-bound interval of $f\in F$,
$X$ the other columns, $\ell_F=\sum_f\min(B_{i_1f}\underline w_f,B_{i_1f}\overline w_f)$,
$h_F=\sum_f\max(B_{i_1f}\underline w_f,B_{i_1f}\overline w_f)$ and
$\overline{A_ip}=\max_{p\in\mathcal P_h}A_ip=\sum_j\max(A_{ij}\underline p_j,A_{ij}\overline p_j)$. Define the
\emph{static robust LP}
\begin{equation}
\begin{aligned}
\varphi_R=\max_{0\le v\le\overline u_X}\ d_{2,X}^\top v\quad\text{s.t.}\quad
&B_{iX}v\le b_i-\overline{A_ip}\qquad (i\notin\{i_1,i_2\},\ B_{iF}=0),\\
&B_{i_1X}v\le b_{i_1}-\overline{A_{i_1}p}-\ell_F,\qquad
B_{i_2X}v\le b_{i_2}-\overline{A_{i_2}p}+h_F .
\end{aligned}
\label{eq:phiR}
\end{equation}

\begin{lemma}[Robust value-function bound]\label{lem:robust}
For every $p\in\mathcal P_h$: $\max\{d_2^\top u:u\in\mathcal U(p)\}\ge\varphi_R$.
\end{lemma}
\begin{proof}
If (\ref{eq:phiR}) is infeasible, $\varphi_R=-\infty$ and there is nothing to prove. Let $v$ be
optimal in (\ref{eq:phiR}) and fix $p$. Rows without $F$: $A_ip+B_{iX}v\le\overline{A_ip}+B_{iX}v\le b_i$.
Equality pair: write $g=A_{i_1}p+B_{i_1X}v$, so the pair reads
$-b_{i_2}\le g+B_{i_1F}w\le b_{i_1}$. As $w$ ranges over the box $\prod_f[\underline w_f,\overline w_f]$,
$B_{i_1F}w$ covers the whole interval $[\ell_F,h_F]$ (a linear image of a box is an interval).
The constraints of (\ref{eq:phiR}) give $g\le b_{i_1}-\ell_F$ (row $i_1$) and
$-g\le b_{i_2}+h_F$ (row $i_2$, since $A_{i_2}=-A_{i_1}$, $B_{i_2X}=-B_{i_1X}$), i.e.\
$[-b_{i_2}-g,\ b_{i_1}-g]\cap[\ell_F,h_F]\neq\emptyset$ because $b_{i_1}+b_{i_2}\ge0$. Pick $w(p)$
in the box with $B_{i_1F}w(p)$ in this intersection; own-bound rows of $F$ hold by
construction. Hence $u=(v,w(p))\in\mathcal U(p)$ and, as $d_{2,F}=0$,
$\max_{u'\in\mathcal U(p)}d_2^\top u'\ge d_{2,X}^\top v=\varphi_R$.
\end{proof}
The dual multipliers $\mu$ of (\ref{eq:phiR}) give the vertex
$p^{\rm adv}_j=\overline p_j$ if $(A^\top\mu)_j>0$, else $\underline p_j$, which maximises $\mu^\top Ap$
over the box; $F(p^{\rm adv})$ is the first incumbent (a heuristic; any feasible $p$ is valid).

\subsection{Validity of the master bound}
\begin{theorem}[Master relaxation]\label{thm:relax}
Let $(p,u^\ast)$ be bilevel feasible, i.e.\ $p\in\mathcal P_h$ and $u^\ast$ optimal for the
follower at $p$, and let all cuts in $\mathcal C$ be generated as above. Then there is
$z$ such that $(p,u^\ast,z)$ is feasible for $\mathrm{MP}(\mathcal C)$. Consequently
$\mathrm{UB}(\mathcal C):=\max\mathrm{MP}(\mathcal C)\ge\mathrm{OPT}$.
\end{theorem}
\begin{proof}
Let $\lambda^\ast$ satisfy (\ref{eq:kkt}) with $u^\ast$, supported on $\mathcal R$. Set
$z^u_k=\mathbb 1[u^\ast_k>0]$ and $z^\lambda_i=\mathbb 1[\lambda^\ast_i>0]$.
(\ref{eq:M1}) holds by feasibility. (\ref{eq:M2}): if $u^\ast_k>0$ then $z^u_k=1$ and
$u^\ast_k\le\overline u_k$ by (\ref{eq:genb}); otherwise $u_k^\ast=0$.
(\ref{eq:M3}): if $\lambda^\ast_i>0$, complementarity gives zero slack, which is the
requirement for $z^\lambda_i=1$; if $z^\lambda_i=0$ the slack is at most $S_i$ by (\ref{eq:S}).
(\ref{eq:M4}): if $\lambda^\ast_a,\lambda^\ast_j>0$ both rows are binding,
$[A_a\,B_a](p;u^\ast)=b_a$ and $-[A_a\,B_a](p;u^\ast)=b_j$, so $b_a+b_j=0$, contradicting
$(a,j)\in\Pi$. (\ref{eq:M5}): $d_2^\top u^\ast=\max_{\mathcal U(p)}d_2^\top u\ge\varphi_R$ by
Lemma~\ref{lem:robust}. (\ref{eq:M6}): suppose a cut $(C_u,C_\ell)$ is violated, i.e.\
$z^u_k=1$ for all $k\in C_u$ and $z^\lambda_i=0$ for all $i\in C_\ell$. Then $u^\ast_k>0$ on $C_u$,
so $(B^\top\lambda^\ast)_k=d_{2,k}$ there by complementarity, and $\lambda^\ast_i=0$ on $C_\ell$;
together with $\lambda^\ast\ge0$, $B^\top\lambda^\ast\ge d_2$ this means $\lambda^\ast$ solves
$\mathrm{DF}(C_u,C_\ell)$, contradicting how the cut was generated. Hence $(p,u^\ast,z)$ is
master-feasible and its objective, the leader value of $(p,u^\ast)$, is at most
$\mathrm{UB}(\mathcal C)$. Taking $(p,u^\ast)$ optimal for (\ref{eq:bilevel}) gives the claim.
\end{proof}

\begin{lemma}[Attained incumbent]\label{lem:attained}
For every $p\in\mathcal P_h$, $F(p)\le\mathrm{OPT}$ and $F(p)$ is attained by the bilevel
feasible point $(p,u(p))$ returned by \texttt{follower\_response}.
\end{lemma}
\begin{proof}
$u(p)$ maximises $d_2^\top u$ over $\mathcal U(p)$ and, among those maximisers, $d_1^\top u$;
so $(p,u(p))$ is feasible for (\ref{eq:bilevel}) and its value $F(p)$ is at most the
optimum.
\end{proof}

\begin{theorem}[Certificate and finite termination of phase 1]\label{thm:benders}
Let $\mathrm{LB}_t=\max_{\tau\le t}F(\hat p_\tau)$ (incumbents, also local search) and
$\mathrm{UB}_t=\min_{\tau\le t}\mathrm{UB}(\mathcal C_\tau)$. Then for every iteration $t$,
\begin{equation}
\mathrm{LB}_t\le\mathrm{OPT}\le\mathrm{UB}_t,
\label{eq:sandwich}
\end{equation}
so $(\mathrm{UB}_t-\mathrm{LB}_t)/\max(1,|\mathrm{UB}_t|)$ is a certified relative gap.
In exact arithmetic and without the lift to phase~2, phase~1 terminates after at most
$2^{n_u+|\mathcal R|}$ iterations, either
(a) because $\mathrm{DF}(J_u(\hat u),J_\ell(\hat s))$ is feasible, in which case $(\hat p,\hat u)$ is
optimal and $\mathrm{LB}=\mathrm{UB}=\mathrm{OPT}$, or (b) because the gap is below the tolerance.
\end{theorem}
\begin{proof}
(\ref{eq:sandwich}) follows from Lemma~\ref{lem:attained} and Theorem~\ref{thm:relax}
(the cut set only grows, so every $\mathrm{UB}(\mathcal C_\tau)$ is a valid bound).

(a) If $\mathrm{DF}(J_u(\hat u),J_\ell(\hat s))$ has a solution $\lambda$, then $\lambda\ge0$,
$B^\top\lambda\ge d_2$, $\lambda_i=0$ wherever $\hat s_i>0$ and $(B^\top\lambda-d_2)_k=0$ wherever
$\hat u_k>0$: (\ref{eq:kkt}) holds, so $\hat u$ is follower-optimal at $\hat p$ and
$(\hat p,\hat u)$ is bilevel feasible with value $\mathrm{UB}(\mathcal C_t)$. With
(\ref{eq:sandwich}), $\mathrm{UB}(\mathcal C_t)=\mathrm{OPT}$.

Otherwise a cut is added that is violated by the current binary vector $\hat z$:
consider $J^z_u=\{k:\hat z^u_k=1\}\supseteq J_u(\hat u)$ (by (\ref{eq:M2})) and
$J^z_\ell=\{i:\hat z^\lambda_i=0\}\supseteq J_\ell(\hat s)$ (by (\ref{eq:M3})). Restricting a
system by fewer equalities/zeros is a relaxation, so $\mathrm{DF}(J^z_u,J^z_\ell)$ feasible
would imply $\mathrm{DF}(J_u(\hat u),J_\ell(\hat s))$ feasible; hence $\mathrm{DF}(J^z_u,J^z_\ell)$ is
infeasible. The implementation uses the Farkas subsets of $(J_u(\hat u),J_\ell(\hat s))$ if
$\hat z$ violates the resulting cut and otherwise those of $(J^z_u,J^z_\ell)$; in both cases
$C_u\subseteq J^z_u$ and $C_\ell\subseteq J^z_\ell$, so the left-hand side of (\ref{eq:M6}) is $0$
at $\hat z$: the cut removes $\hat z$. By Theorem~\ref{thm:relax} it never removes the binary
vector of an optimal KKT pair, so the master stays feasible. Each iteration removes at
least one of the finitely many binary vectors, which proves finite termination.
\end{proof}
\begin{remark}
$\mathrm{UB}$ is taken from the master solution with relative MIP gap $10^{-8}$; the bound is
therefore valid up to that tolerance. Tolerances $\hat u_k>10^{-5}\overline u_k$ and
$\hat s_i>10^{-5}S_i$ define $J_u,J_\ell$ numerically.
\end{remark}

\subsection{Phase 2: dual lifting}
If (\ref{eq:phiR}) is infeasible or phase~1 exceeds 15 iterations, the master is lifted
with dual variables $\lambda\in\mathbb R^{\mathcal R}_{\ge0}$ and the constraints
\begin{equation}
B_{\mathcal R}^\top\lambda\ge d_2,\qquad B_{\mathcal R}^\top\lambda-d_2\le L_d(\mathbf 1-z^u),\qquad
\lambda\le L_d\,z^\lambda,
\label{eq:lift}
\end{equation}
($L_d=5\cdot10^4$; a big-M on the dual side only), keeping (\ref{eq:M1})--(\ref{eq:M6}).

\begin{lemma}[A uniform dual bound exists]\label{lem:dualbound}
The dual feasible set $\mathcal D=\{\lambda\in\mathbb R^{\mathcal R}_{\ge0}:B_{\mathcal R}^\top\lambda\ge d_2\}$
does not depend on $p$. For every $p$ there is an optimal dual solution at a vertex of
$\mathcal D$. Hence $L^\ast=\max_{\lambda\in\mathrm{vert}(\mathcal D)}\max(\|\lambda\|_\infty,\|B^\top\lambda-d_2\|_\infty)<\infty$
bounds an optimal KKT multiplier for every $p$ simultaneously.
\end{lemma}
\begin{proof}
$\mathcal D$ involves only $B$ and $d_2$. It lies in the non-negative orthant, so it is a
pointed polyhedron with finitely many vertices. Under Assumption~\ref{as:feas} the
follower LP is feasible and bounded (by (\ref{eq:genb})), so by strong duality the dual
$\min_{\lambda\in\mathcal D}\lambda^\top(b-Ap)$ has a finite optimum, attained at a vertex
(fundamental theorem of LP).
\end{proof}
\begin{theorem}[Phase-2 exactness]\label{thm:phase2}
Every feasible point of the lifted master is bilevel feasible. If $L_d\ge L^\ast$, the
lifted master optimum equals $\mathrm{OPT}$.
\end{theorem}
\begin{proof}
At a lifted-master point: $z^u_k=0\Rightarrow u_k=0$ by (\ref{eq:M2}); $z^u_k=1\Rightarrow
(B^\top\lambda-d_2)_k=0$ by (\ref{eq:lift}); so $u_k(B^\top\lambda-d_2)_k=0$. Likewise
$z^\lambda_i=0\Rightarrow\lambda_i=0$ and $z^\lambda_i=1\Rightarrow$ zero slack by (\ref{eq:M3}). Hence (\ref{eq:kkt})
holds and $u$ is follower-optimal: the point is bilevel feasible and its value is
$\le\mathrm{OPT}$. Conversely take an optimal $(p,u^\ast)$ and a vertex multiplier $\lambda^\ast$
(Lemma~\ref{lem:dualbound}); with $z$ as in the proof of Theorem~\ref{thm:relax},
(\ref{eq:M1})--(\ref{eq:M6}) hold and (\ref{eq:lift}) holds because
$\lambda^\ast_i\le L^\ast\le L_d$ and $(B^\top\lambda^\ast-d_2)_k\le L^\ast\le L_d$ (and $=0$ when $z^u_k=1$).
So the lifted master value is $\ge\mathrm{OPT}$.
\end{proof}
Phase 2 was needed in 1 of 48 problems (hour 21, upper); its certificate is conditional on
$L_d\ge L^\ast$, whereas phase 1 is unconditional.

\subsection{Algorithm}
\begin{enumerate}
\item Compute $S$ (\ref{eq:S}), $\Pi$, $\varphi_R$ and $p^{\rm adv}$ (\ref{eq:phiR}); $\mathrm{LB}\gets F(p^{\rm adv})$, $\mathcal C\gets\emptyset$.
\item Solve $\mathrm{MP}(\mathcal C)$ (lifted if phase 2) $\to(\hat p,\hat u,\hat z)$, $\mathrm{UB}\gets\min(\mathrm{UB},\mathrm{val})$.
\item $\mathrm{LB}\gets\max(\mathrm{LB},F(\hat p))$; in iteration 1, one-flip local search over vertices.
\item Stop if phase 2, if the gap $\le10^{-3}$, or if $\mathrm{DF}(J_u(\hat u),J_\ell(\hat s))$ is feasible.
\item Otherwise add the Farkas cut (\ref{eq:M6}) and go to 2; lift (phase 2) after 15 iterations.
\end{enumerate}

\section{AC Model, Security and the AC-Correction Loop}
\label{app:ac}

\subsection{AC power flow}
[gridfm-datakit + lightsim2grid] For an operating point $(h,s,p,u)$ the bus voltages
$V_i\angle\theta_i$ solve, with the topology-$s$ admittance matrix $Y^s=G^s+jB^s$,
\begin{align}
P^{\rm inj}_i&=V_i\textstyle\sum_{j}V_j\big(G^s_{ij}\cos\theta_{ij}+B^s_{ij}\sin\theta_{ij}\big),\label{eq:acp}\\
Q^{\rm inj}_i&=V_i\textstyle\sum_{j}V_j\big(G^s_{ij}\sin\theta_{ij}-B^s_{ij}\cos\theta_{ij}\big),\label{eq:acq}
\end{align}
where $P^{\rm inj}=C_Gu+C_RP_R-C_DP_D$ at PQ and PV buses (slack injection free),
$Q^{\rm inj}=-C_D(\gamma\circ P_D)$ at PQ buses with the fixed load power factor
$\gamma_i=Q^0_{D,i}/P^0_{D,i}$, generator buses at the AC-OPF set-point $V_i=V^{\rm set}_i(h)$, and
$\theta_{n_0}=0$. The contingency doubles the series impedance $r_\ell+jx_\ell$ and halves the
charging susceptance of the corridor (one of two identical circuits lost), consistent with
(\ref{eq:xs}); PV injects at unity power factor.
Branch $\ell=(i,j)$ in the $\pi$-model carries
$S^{\rm f}_\ell=V_i\angle\theta_i\,(I^{\rm f}_\ell)^\ast$, $S^{\rm t}_\ell=V_j\angle\theta_j\,(I^{\rm t}_\ell)^\ast$ and
$S_\ell=\max(|S^{\rm f}_\ell|,|S^{\rm t}_\ell|)$.

\subsection{Security quantities}
[\texttt{acrss.read\_ac}, \texttt{genco.evaluate}] With PQ buses $\mathcal N_{PQ}$, limits
$[\underline V_i,\overline V_i]$ and tolerance $\tau_V=0.01$\,p.u.,
\begin{equation}
q_{\rm load}=\max_{\ell}\frac{S_\ell}{\overline S_\ell},\qquad
q_{\rm vlow}=\max_{i\in\mathcal N_{PQ}}(\underline V_i-V_i),\qquad
q_{\rm vhigh}=\max_{i\in\mathcal N_{PQ}}(V_i-\overline V_i),
\label{eq:q}
\end{equation}
and the scenario is \emph{secure} iff
\begin{equation}
q_{\rm load}\le1,\qquad q_{\rm vlow}\le\tau_V,\qquad q_{\rm vhigh}\le\tau_V .
\label{eq:secure}
\end{equation}
For verified scenarios the slack output is also checked, $P_{g_0}\le P^{\max}_{g_0}$.

\subsection{AC-correction update}
[\texttt{acrss.ac\_rss}] At iteration $t$ the DC model uses $\kappa^t$ and $\overline u^t_{g_0}$.
For each verified overload ($q_{{\rm load},\ell}=S_\ell/\overline S_\ell>1$ in topology $s$),
with $\varepsilon=0.01$ and floor $\kappa_{\min}=0.8$,
\begin{equation}
\kappa^{t+1}_{s\ell}=\max\Big(\kappa_{\min},\ \min\Big(\kappa^t_{s\ell},\ \kappa^t_{s\ell}\,\frac{1-\varepsilon}{S_\ell/\overline S_\ell}\Big)\Big),
\qquad
\overline u^{t+1}_{g_0}=\overline u^{t}_{g_0}-(1+\varepsilon)\max_{\rm verified}(P_{g_0}-P^{\max}_{g_0})^+ .
\label{eq:derate}
\end{equation}

\begin{proposition}[Nesting]\label{prop:nest}
If $\kappa'\le\kappa$ componentwise and $\overline u'\le\overline u$, then
$\mathcal U'(p)\subseteq\mathcal U(p)$ for all $p$, $\varphi'\le\varphi$, $\psi'\ge\psi$, and
$[\underline\alpha'_h,\overline\alpha'_h]\subseteq[\underline\alpha_h,\overline\alpha_h]$.
In particular the iterates of (\ref{eq:derate}) produce nested intervals, all inside the
DC interval.
\end{proposition}
\begin{proof}
Only right-hand sides change: the line rows (\ref{eq:line}) and the bound $\overline u_{g_0}$
become tighter, so every $u\in\mathcal U'(p)$ satisfies the constraints of $\mathcal U(p)$.
A maximum over a subset is not larger and a minimum over a subset is not smaller, hence
$\varphi'\le\varphi$, $\psi'\ge\psi$, and by (\ref{eq:rssdef}) $\overline\alpha'\le\overline\alpha$,
$\underline\alpha'\ge\underline\alpha$. The update (\ref{eq:derate}) is non-increasing in both
$\kappa$ and $\overline u_{g_0}$.
\end{proof}
\begin{corollary}[Emptiness is robust to merging the cuts]\label{cor:empty}
The two limits are corrected with their own factors $\kappa^{\rm up},\kappa^{\rm lo}$. If
$\underline\alpha_h(\kappa^{\rm lo})>\overline\alpha_h(\kappa^{\rm up})$, then with the merged factors
$\kappa=\min(\kappa^{\rm up},\kappa^{\rm lo})$ the robust secure set is also empty:
$\underline\alpha_h(\kappa)\ge\underline\alpha_h(\kappa^{\rm lo})>\overline\alpha_h(\kappa^{\rm up})\ge\overline\alpha_h(\kappa)$
by Proposition~\ref{prop:nest}, and Theorem~\ref{thm:rss} applies.
\end{corollary}
\begin{proposition}[Termination and what is certified]\label{prop:loop}
Each $\kappa_{s\ell}$ can strictly decrease at most
$\lceil\ln\kappa_{\min}/\ln(1-\varepsilon)\rceil=23$
times before it reaches the floor, and the
loop is capped at 6 iterations. If it stops at iteration $t$ because no verified scenario
violates, then: (i) $\alpha^t$ is the certified optimum of the corrected DC bilevel
(Theorem~\ref{thm:benders}); (ii) every flagged pool scenario is AC-secure, verified by exact
power flow; (iii) every cleared pool scenario that is exchangeable with the calibration
data is insecure with probability at most $\alpha$ (Theorem~\ref{thm:conf}), and the expected
number of insecure cleared scenarios in a pool of $N$ is at most $\alpha N$; (iv) with the
audit enabled (all our experiments) every pool scenario is solved exactly, and a zero
audit count means the whole final pool is thermally secure.
\end{proposition}
\begin{proof}
A strict decrease multiplies $\kappa_{s\ell}$ by $(1-\varepsilon)/q<1-\varepsilon$ since $q>1$; after
$n$ decreases $\kappa_{s\ell}\le(1-\varepsilon)^n$, which is below $\kappa_{\min}$ for
$n\ge\ln\kappa_{\min}/\ln(1-\varepsilon)$. (i) is Theorem~\ref{thm:benders} applied to the data
of iteration $t$; (ii) is the stopping rule; (iii) is Theorem~\ref{thm:conf} plus linearity
of expectation; (iv) is the definition of the audit.
\end{proof}
\begin{remark}
(\ref{eq:derate}) is a correction heuristic, not an exact AC robust counterpart: it
tightens the \emph{active-power} limit of a branch whose \emph{apparent} power is too high.
If the overload is mostly reactive, the floor $\kappa_{\min}$ is reached (it was reached in
every lower-limit problem of E4), and the corrected limits are conservative.
\end{remark}

\section{GENCO Trajectory, Recalibration and Conformal Screening}
\label{app:conf}

\subsection{Corrective-layer trajectory}
GENCO (\texttt{GNS\_heterogeneous}) refines a state estimate $\hat x^{(l)}=(\hat V^{(l)},\hat\theta^{(l)})$ over
$K$ layers. In layer $l$ it evaluates the AC mismatch of (\ref{eq:acp})--(\ref{eq:acq}),
$\Delta P^{(l)}_i=P^{\rm inj}_i-P_i(\hat x^{(l)})$, $\Delta Q^{(l)}_i=Q^{\rm inj}_i-Q_i(\hat x^{(l)})$, and feeds it back
into the bus embeddings. A forward hook on its residual layer records
\begin{equation}
r^{(l)}_i=\sqrt{(\Delta P^{(l)}_i)^2+(\Delta Q^{(l)}_i)^2},\qquad
\bar r^{(l)}=\frac1{|\mathcal N|}\sum_i r^{(l)}_i,\qquad \rho=\bar r^{(K)}/\bar r^{(1)}.
\label{eq:traj}
\end{equation}
GENCO's predictions $\hat q_k$ of (\ref{eq:q}) use its final state and graphkit's branch-flow
operator.

\subsection{Physics-trajectory recalibration (PTR)}
Let $z=(\hat q_{\rm load},\hat q_{\rm vlow},\hat q_{\rm vhigh},\hat q_{\rm qex},\hat q_{\rm slack},\bar r^{(1)},\dots,\bar r^{(K)},\rho)$
be GENCO-internal features (no hour, no topology label). On a \emph{tune} set
$\mathcal T$ solved exactly, fit for each $k$ a gradient-boosted regression ensemble
\begin{equation}
g_k=\arg\min_{g\in\mathcal F}\sum_{j\in\mathcal T}\big(q_{k,j}-\hat q_{k,j}-g(z_j)\big)^2
\quad(300\text{ depth-3 trees, rate }0.05,\text{ subsample }0.8),
\label{eq:ptr}
\end{equation}
and set $\tilde q_k=\hat q_k+g_k(z)$. Plain conformal uses $g_k\equiv0$; the residual-normalised
variant uses $g_k\equiv0$ with a scale $\sigma_k+\beta_k\bar r^{(K)}$ tuned on $\mathcal T$.

\subsection{Split-conformal screening}
Let $\ell_{\rm load}=1$, $\ell_{\rm vlow}=\ell_{\rm vhigh}=\tau_V$, a level $\alpha\in(0,1)$,
$a_k=\alpha/3$, and a calibration set $\mathcal C$ of $n$ scenarios solved exactly,
independent of $\mathcal T$. With scale $\sigma_k(z)>0$ ($\equiv1$ for plain and PTR),
\begin{equation}
s_{k,j}=\frac{q_{k,j}-\tilde q_{k,j}}{\sigma_k(z_j)},\qquad
c_k=\text{the }\big\lceil(n+1)(1-a_k)\big\rceil\text{-th smallest of }\{s_{k,j}\}_{j\in\mathcal C}\ (+\infty\text{ if }>n),
\label{eq:ck}
\end{equation}
and a new scenario is \emph{flagged} (sent to exact AC) iff
\begin{equation}
\exists k:\ \tilde q_k+c_k\,\sigma_k(z)>\ell_k ;
\label{eq:flag}
\end{equation}
otherwise it is \emph{cleared}.

\begin{lemma}[Split-conformal quantile]\label{lem:conf}
Let $S_1,\dots,S_{n+1}$ be exchangeable real random variables, $k=\lceil(n+1)(1-a)\rceil$ and
$c$ the $k$-th smallest of $S_1,\dots,S_n$ ($c=+\infty$ if $k>n$). Then
$\Pr[S_{n+1}>c]\le a$.
\end{lemma}
\begin{proof}
If $k>n$ the event is empty. Otherwise, if $S_{n+1}>c$ then at least $k$ of
$S_1,\dots,S_n$ are strictly smaller than $S_{n+1}$, i.e.\ $R_{n+1}\ge k$ with
$R_i=\#\{j\le n+1:S_j<S_i\}$. Let $v_{(1)}\le\dots\le v_{(n+1)}$ be the ordered values. If
$R_i\ge k$ then $S_i>v_{(k)}$, which is possible only for indices whose value sits in
positions $k+1,\dots,n+1$; hence $\sum_{i=1}^{n+1}\mathbb 1[R_i\ge k]\le n+1-k$ surely. By
exchangeability $\Pr[R_i\ge k]$ is the same for all $i$, so
$(n+1)\Pr[R_{n+1}\ge k]=\mathbb E\sum_i\mathbb 1[R_i\ge k]\le n+1-k$ and
$\Pr[S_{n+1}>c]\le\Pr[R_{n+1}\ge k]\le\frac{n+1-k}{n+1}\le a$, because $k\ge(n+1)(1-a)$.
\end{proof}

\begin{theorem}[No-miss guarantee]\label{thm:conf}
Let the calibration scenarios and a test scenario be exchangeable (e.g.\ i.i.d.\ draws of
the screening distribution) and let $g_k,\sigma_k$ be fixed given $\mathcal T$ (fitted
without $\mathcal C$ and the test scenario). Then
\begin{equation}
\Pr[\text{test scenario insecure by (\ref{eq:secure}) and cleared}]\le\alpha .
\label{eq:guarantee}
\end{equation}
\end{theorem}
\begin{proof}
Condition on $\mathcal T$. Then $g_k$ and $\sigma_k$ are fixed functions, so the scores
$s_{k,1},\dots,s_{k,n},s_{k,n+1}$ (the last for the test scenario) are exchangeable for each $k$.
If the test scenario is insecure, $q_k>\ell_k$ for some $k$; if it is cleared,
$\tilde q_k+c_k\sigma_k\le\ell_k<q_k$, i.e.\ $s_{k,n+1}>c_k$. Hence
\[
\{\text{insecure and cleared}\}\subseteq\bigcup_{k}\{s_{k,n+1}>c_k\},\qquad
\Pr[\cdot\mid\mathcal T]\le\sum_k\Pr[s_{k,n+1}>c_k\mid\mathcal T]\le\sum_k a_k=\alpha
\]
by the union bound and Lemma~\ref{lem:conf}. Taking the expectation over $\mathcal T$
gives (\ref{eq:guarantee}).
\end{proof}
\begin{corollary}\label{cor:any}
The guarantee holds for \emph{every} recalibrator $g_k$ (including $g_k\equiv0$ and all
ablation variants), because the proof uses only that $g_k$ is fixed before calibration.
A better $g_k$ cannot weaken it; it only reduces the flag rate. For $N$ exchangeable
scenarios the expected number of insecure-and-cleared ones is at most $\alpha N$.
\end{corollary}
\begin{remark}
The guarantee is marginal (over the screening distribution), not conditional on the
hour or topology. For the unseen topology the exchangeability assumption does not hold;
the measured miss rates there are an empirical robustness check.
\end{remark}

\section{GENCO-Gradient Adversary}
\label{app:adv}
[\texttt{genco.adversarial}] With the dispatch fixed, let $\hat\ell_e(p)$, $e\in\mathcal E$, be the
GENCO-predicted loadings $S_\ell/\overline S_\ell$ of all branches in the three topologies
(the reactive load moves with $P_D$ at fixed power factor). The objective is the smooth
maximum
\begin{equation}
J(p)=\tau\log\sum_{e\in\mathcal E}\exp\big(\hat\ell_e(p)/\tau\big),\qquad\tau=0.02 .
\label{eq:J}
\end{equation}
\begin{lemma}[Smooth maximum]\label{lem:lse}
$\max_e\hat\ell_e\le J\le\max_e\hat\ell_e+\tau\log|\mathcal E|$, and
$\nabla_pJ=\sum_e\pi_e\nabla_p\hat\ell_e$ with $\pi=\mathrm{softmax}(\hat\ell/\tau)$ in the simplex.
\end{lemma}
\begin{proof}
With $m=\max_e\hat\ell_e$: $e^{m/\tau}\le\sum_e e^{\hat\ell_e/\tau}\le|\mathcal E|e^{m/\tau}$; take
$\tau\log$. The gradient follows from the chain rule,
$\partial J/\partial\hat\ell_e=e^{\hat\ell_e/\tau}/\sum_{e'}e^{\hat\ell_{e'}/\tau}=\pi_e$.
\end{proof}
The update is projected sign-gradient ascent on the box,
\begin{equation}
p\leftarrow\Pi_{\mathcal P_h}\big(p+\eta\,(\overline p-\underline p)\circ\mathrm{sign}\,\nabla_pJ(p)\big),\qquad
\Pi_{\mathcal P_h}(y)=\min(\max(y,\underline p),\overline p),\ \eta=0.25,
\label{eq:adv}
\end{equation}
for 6 steps from the riskiest pool points; the result is rounded to the nearest vertex,
$p_j\leftarrow\overline p_j$ if $p_j>\tfrac12(\underline p_j+\overline p_j)$ else $\underline p_j$ (by
Lemma~\ref{lem:vertex} the DC worst case is a vertex), the follower response $u(p)$ is
recomputed and the point is verified by exact AC. $\Pi_{\mathcal P_h}$ is the Euclidean
projection onto a box (separable). The adversary is a heuristic; nothing is certified from
it without exact verification.

\section{Second-Life EV Batteries}
\label{app:bess}
[\texttt{bess.bess\_power}, \texttt{rss.\_add\_bess}] At each PV bus $j\in\mathcal R$ a battery of
power $P^{\rm b}_j=\mathrm{SoH}_j\cdot0.35\cdot\bar P_{R,j}$ adds discharge and charge variables
$0\le u^{\rm dis}_j,u^{\rm ch}_j\le P^{\rm b}_j$ to the follower. Its net injection
$u^{\rm dis}-u^{\rm ch}$ enters (\ref{eq:flow}) through $C_R$ and the balance (\ref{eq:bal});
it does not enter the index, $d_{\rm dis}=d_{\rm ch}=0$.
\begin{proposition}[Batteries can only widen the RSS]\label{prop:bess}
$\overline\alpha^{\rm b}_h\ge\overline\alpha_h$ and $\underline\alpha^{\rm b}_h\le\underline\alpha_h$.
\end{proposition}
\begin{proof}
For every $u\in\mathcal U(p)$, $(u,0,0)$ is feasible with batteries (zero battery
variables leave all rows unchanged), with the same index. Hence the index range at each
$p$ can only grow, $\varphi^{\rm b}\ge\varphi$ and $\psi^{\rm b}\le\psi$, and (\ref{eq:rssdef})
gives the claim.
\end{proof}
